% Part: lambda-calculus
% Chapter: lambda-definability
% Section: arithmetical-functions

\documentclass[../../../include/open-logic-section]{subfiles}

\begin{document}

\olfileid[fr]{lam}{rep}{arf}
\olsection{Fonctions arithmétiques \texorpdfstring{$\lambda$}{λ}-définissables}

\begin{prop}
\ollabel{prop:succ-ld}
La fonction successeur~$\Succ$ est
$\lambda$-définissable.
\end{prop}

\begin{proof}
Le terme suivant $\lambda$-définit
la fonction successeur :
\begin{align*}
  \fn{Succ} & \ident
    \lambd[a][\lambd[fx][f ({a} f x)]].
  \intertext{Avec nos conventions, cette
    expression abrège :}
  \fn{Succ} & \ident
    \lambd[a][\lambd[f][\lambd[x][(f (({a} f) x))]]].
  \intertext{Le terme $\fn{Succ}$ reçoit un
    nombre~$a$ et renvoie la fonction
    $\lambd[fx][f ({a} f x)]$. Cette fonction
    n'est pas elle-même un numéral de Church.
    Cependant, si $a$ est un numéral de Church,
    l'application se réduit à un tel numéral :}
  (\lambd[a][\lambd[fx][f ({a} f x)]])\,\num{n}
    & \redone \lambd[fx][f ({\num{n}} f x)].
  \intertext{Le sous-terme $\num{n}fx$ se
    réduit à $f^n(x)$ en deux contractions
    à la racine, puisque
    $\num{n}=\lambd[fx][f^n(x)]$. Ainsi :}
  \fn{Succ}\,\num{n}
    & \red \lambd[fx][f(f^n(x))],
\end{align*}
c'est-à-dire à~$\num{n+1}$.
\end{proof}

\begin{ex}
Calculons $\fn{Succ}$ appliqué à~$\num{0}$,
c'est-à-dire à $\lambd[fx][x]$, en écrivant
les termes sans abréviation :
\begin{align*}
  \fn{Succ}\,\num{0}
    & \ident
    (\lambd[a][\lambd[f][\lambd[x][(f (({a} f) x))]]])
    (\lambd[f][\lambd[x][x]]) \\
  & \redone
    \lambd[f][\lambd[x][(f (({(\lambd[f][\lambd[x][x]])} f) x))]] \\
  & \redone
    \lambd[f][\lambd[x][(f ({(\lambd[x][x])} x))]] \\
  & \redone \lambd[f][\lambd[x][(f x)]]
    \ident \num{1}.
\end{align*}
\end{ex}

\begin{prob}
Le terme
\begin{align*}
  \fn{Succ}' & \ident
    \lambd[{n}][\lambd[fx][{n} f (f x)]]
\end{align*}
$\lambda$-définit aussi la fonction
successeur. Expliquer pourquoi.
\end{prob}

\begin{prop}
\ollabel{prop:add-ld}
La fonction addition~$\Add$ est
$\lambda$-définissable.
\end{prop}

\begin{proof}
L'addition est $\lambda$-définie
par l'un ou l'autre des termes :
\begin{align*}
  \fn{Add} & \ident
    \lambd[{a}{b}][\lambd[fx][{a} f ({b} f x)]]
\intertext{ou bien :}
  \fn{Add}' & \ident
    \lambd[{a}{b}][{a}\,\fn{Succ}\,{b}].
\intertext{Le premier terme reçoit deux
  numéraux $a$ et~$b$, puis renvoie une
  fonction qui reçoit $f$ et~$x$ et produit
  $af(bfx)$. Pour $a=\num{n}$ et
  $b=\num{m}$, cela donne
  $f^n(f^m(x))=f^{n+m}(x)$. En détail :}
  (\lambd[{a}{b}][\lambd[fx][{a} f ({b} f x)]])
    \num{n}\,\num{m}
    & \red \lambd[fx][\num{n}\,f(\num{m}\,f x)] \\
  & \red \lambd[fx][\num{n}\,f(f^m x)] \\
  & \red \lambd[fx][f^n(f^m x)]
    \ident \num{n+m}.
\intertext{La seconde représentation
  fonctionne autrement. Appliquée à
  $\num{n}$ et $\num{m}$, elle donne :}
  \fn{Add}'\,\num{n}\,\num{m}
    & \red \num{n}\,\fn{Succ}\,\num{m}.
\intertext{Or $\num{n}f x$ se réduit à
  $f^n(x)$. Par conséquent :}
  \num{n}\,\fn{Succ}\,\num{m}
    & \red \fn{Succ}^n(\num{m}).
\end{align*}
Comme $\fn{Succ}$ $\lambda$-définit
la fonction successeur, $n$ itérations
à partir de~$m$ conduisent à~$n+m$;
le terme se réduit donc à~$\num{n+m}$.
\end{proof}

\begin{prop}
\ollabel{prop:mult-ld}
La multiplication est $\lambda$-définie
par le terme
\[
  \fn{Mult} \ident
  \lambd[ab][\lambd[fx][a (b f) x]].
\]
\end{prop}

\begin{proof}
Appliquons $\fn{Mult}$ aux numéraux
$\num{n}$ et $\num{m}$. On obtient
$\lambd[fx][\num{n}(\num{m}\,f)x]$.
Le terme $\num{m}f$ est une fonction
qui applique~$f$ $m$ fois à son argument;
$\num{n}(\num{m}f)x$ itère cette fonction
$n$ fois à partir de~$x$. Il applique
donc~$f$ à~$x$ $n\cdot m$ fois, et la
forme normale obtenue est $\num{nm}$.
\end{proof}

Le terme définissant la multiplication
admet une contraction~$\eta$, puisque $x$
n'est pas libre dans $a(bf)$ :
\[
  \fn{Mult} \eredone
  \lambd[ab][\lambd[f][a (b f)]].
\]

\begin{prob}
La multiplication peut aussi être
$\lambda$-définie par
\[
  \fn{Mult}' \ident
  \lambd[ab][a (\fn{Add}\, b) \num{0}].
\]
Expliquer pourquoi ce terme convient.
\end{prob}

Une vérification de l'arité consiste à comparer
$\fn{Mult}'\,\num{2}\,\num{3}$ et
$\fn{Mult}'\,\num{2}\,\num{4}$.
Le deuxième argument fixe l'addend et le premier
fixe le nombre d'additions.

La définition de l'exponentiation par
un lambda-terme est étonnamment simple :
\[
  \fn{Exp} \ident \lambd[be][\lambd[fx][e b f x]].
\]
Le premier argument~$b$ est la base, le second~$e$ l'exposant.
Pour $b=\num{n}$, l'application de $b$ à un itérateur $g$
donne un itérateur qui effectue $n$ applications de $g$.
Partons de l'itérateur $f$. Après $k$ applications de $b$,
l'itérateur obtenu effectue $n^k$ applications de $f$.
Au départ, il effectue une application de $f$. À chaque étape,
le nombre d'applications est multiplié par $n$.
Ainsi, avec $e=\num{k}$, le terme $e b f x$ se réduit
à $f^{n^k}(x)$ et $\fn{Exp}\,b\,e$ se réduit au numéral
$\num{n^k}$. On peut aussi définir l'exponentiation
par des multiplications répétées :
\[
  \fn{Exp}' \ident
  \lambd[be][e (\fn{Mult}\, b) \num{1}].
\]
Pour un exposant nul, $\num{0}bfx$ se réduit à $fx$.
On obtient ainsi $\fn{Exp}\,b\,\num{0}\red\num{1}$,
même si $b=\num{0}$. Dans la deuxième définition,
l'itération comporte zéro multiplication et laisse
la valeur initiale $\num{1}$.
Le prédécesseur et la soustraction sont
moins faciles à représenter avec les
numéraux de Church : ils demandent un
codage des paires.

\end{document}
